\documentclass[11pt]{article}
\usepackage{amsmath}
\usepackage{geometry}
\geometry{margin=1in}
\usepackage{xcolor}
\usepackage[colorlinks=true, linkcolor=blue, urlcolor=blue, citecolor=red]{hyperref}
\usepackage[backend=biber,style=numeric,sorting=none,doi=false,url=false]{biblatex}
\addbibresource{references.bib}
\newcommand{\biblinklabel}[1]{%
\iffieldundef{url}
{\iffieldundef{doi}{#1}{\href{https://doi.org/\thefield{doi}}{#1}}}
{\href{\thefield{url}}{#1}}}
% Bibliography labels: red number hyperlinked, brackets applied only once
\DeclareFieldFormat{labelnumber}{\biblinklabel{\textcolor{red}{#1}}}
\DeclareFieldFormat{labelnumberwidth}{\mkbibbrackets{\biblinklabel{\textcolor{red}{#1}}}}
% In-text \cite uses red, bracketed numeric labels linked to bibliography
\DeclareCiteCommand{\cite}
{\usebibmacro{prenote}}
{\mkbibbrackets{\printtext[bibhyperref]{\textcolor{red}{\thefield{labelnumber}}}}}
{\multicitedelim}
{\usebibmacro{postnote}}

\title{MRI Image Quality Metrics (IQMs)}
\author{RadQy Project}
\date{\today}



\begin{document}

\maketitle

\begin{abstract}
This document summarizes commonly used MRI image quality metrics (IQMs) and their mathematical definitions for reproducibility and reporting.
\end{abstract}
\tableofcontents
\bigskip

\section*{Introduction}
\addcontentsline{toc}{section}{Introduction}
This document summarizes commonly used MRI image quality metrics (IQMs) and their mathematical definitions for reproducibility and reporting.

\section{Per-Slice IQMs (Foreground / Background pixels)}
All metrics below accept foreground pixels $\mathbf{F}$ (from segmentation) and background pixels $\mathbf{B}$. NaN/Inf inputs yield NaN. For BG-dependent metrics (SNR1, SNR2, SNR4, CNR, CJV, FBER), if background is missing or unusable (empty, zero variance, non-finite), we fall back to foreground homogeneity $\mu_F/\sigma_F$ to avoid artificial NaNs.

\subsection{Mean (MEAN)}
Foreground mean intensity:
\[
\text{MEAN} = \mu_F = \frac{1}{N_F} \sum_{i} F_i
\]

\subsection{Range (RNG)}
Foreground intensity range:
\[
\text{RNG} = \max(F) - \min(F)
\]

\subsection{Variance (VAR)}
Foreground variance:
\[
\text{VAR} = \sigma_F^2 = \frac{1}{N_F}\sum_i (F_i - \mu_F)^2
\]

\subsection{Coefficient of Variation (CV)}
Foreground coefficient of variation (percent):
\[
\text{CV} = \frac{\sigma_F}{\mu_F} \times 100
\]

\subsection{Contrast per Pixel (CPP)}
Laplacian-filtered contrast, mean response of a $3\times3$ Laplacian:
\[
f_1 = \frac{1}{8}
\begin{bmatrix}
-1 & -1 & -1\\
-1 & 8  & -1\\
-1 & -1 & -1
\end{bmatrix}, \quad
\text{CPP} = \text{mean}\left(f_1 * F\right)
\]
where $*$ is 2D convolution (foreground is median-filtered only when used for PSNR below).

\subsection{Peak Signal-to-Noise Ratio (PSNR)}
\[
\text{PSNR} = 10 \log_{10} \left( \frac{\max(F)^2}{\text{MSE}(F,\hat{F})} \right)
\]
where $\hat{F}$ is either the background $\mathbf{B}$ (if provided) or a median-filtered version of $\mathbf{F}$.

\subsection{SNR1}
Foreground/background std; fallback to foreground stability if background unusable \cite{sadri2020mrqy, kushol2023dsmri}:
\[
\text{SNR1} =
\begin{cases}
\dfrac{\sigma_F}{\sigma_B}, & \sigma_B \text{ finite and } \sigma_B \neq 0\\[6pt]
\dfrac{\mu_F}{\sigma_F}, & \text{otherwise (fg homogeneity)}
\end{cases}
\]

\subsection{SNR2}
Foreground mean over background std; fallback to foreground stability if background unusable:
\[
\text{SNR2} =
\begin{cases}
\dfrac{\mu_F}{\sigma_B}, & \sigma_B \text{ finite and } \sigma_B \neq 0\\[6pt]
\dfrac{\mu_F}{\sigma_F}, & \text{otherwise}
\end{cases}
\]

\subsection{SNR3}
Foreground standard deviation divided by centered-foreground standard deviation (stability check):
\[
\text{SNR3} = \frac{\sigma_F}{\sigma_{F - \mu_F}}
\]

\subsection{SNR4}
Foreground mean over background mean; fallback to foreground stability if background unusable:
\[
\text{SNR4} =
\begin{cases}
\dfrac{\mu_F}{\mu_B}, & \mu_B \text{ finite and } \mu_B \neq 0\\[6pt]
\dfrac{\mu_F}{\sigma_F}, & \text{otherwise}
\end{cases}
\]

\subsection{Contrast-to-Noise Ratio (CNR)}
Foreground/background mean difference over background std; fallback to foreground stability if background unusable:
\[
\text{CNR} =
\begin{cases}
\dfrac{\mu_F - \mu_B}{\sigma_B}, & \sigma_B \text{ finite and } \sigma_B \neq 0\\[6pt]
\dfrac{\mu_F}{\sigma_F}, & \text{otherwise}
\end{cases}
\]

\subsection{Patch Coefficient of Variation (CVP)}
Foreground patch coefficient of variation:
\[
\text{CVP} = \frac{\sigma_F}{\mu_F}
\]

\subsection{Coefficient of Joint Variation (CJV)}
Joint variation between foreground and background; fallback to foreground stability if background unusable:
\[
\text{CJV} =
\begin{cases}
\dfrac{\sigma_F + \sigma_B}{|\mu_F - \mu_B|}, & \text{finite and } |\mu_F - \mu_B| \neq 0\\[6pt]
\dfrac{\mu_F}{\sigma_F}, & \text{otherwise}
\end{cases}
\]

\subsection{Entropy Focus Criterion (EFC)}
Normalized Shannon entropy of foreground:
\[
p_i = \frac{F_i}{\sqrt{\sum_j F_j^2}}, \quad
\text{EFC} = -\frac{1}{\log N_F} \sum_i p_i \log p_i
\]
higher values suggest more ghosting/noise.

\subsection{Foreground-Background Energy Ratio (FBER)}
Median energy ratio; fallback to foreground stability if background unusable:
\[
\text{FBER} =
\begin{cases}
\dfrac{\text{median}(F^2)}{\text{median}(B^2)}, & \text{finite and } \text{median}(B^2)\neq 0\\[6pt]
\dfrac{\mu_F}{\sigma_F}, & \text{otherwise}
\end{cases}
\]

\subsection{Skewness (SKW)}
Foreground skewness (third standardized moment):
\[
\text{SKW} = \frac{\frac{1}{N_F}\sum_i (F_i - \mu_F)^3}{\sigma_F^3}
\]

\subsection{Excess Kurtosis (KURT)}
Foreground excess kurtosis (fourth standardized moment minus 3):
\[
\text{KURT} = \frac{\frac{1}{N_F}\sum_i (F_i - \mu_F)^4}{\sigma_F^4} - 3
\]

\section{Participant-Level QC}
\subsection{Failure Count (FAIL\_FRAC)}
Counts total non-finite IQM evaluations (NaN/Inf) across all slices/metrics for a participant. Higher counts flag poor FG/BG segmentation or corrupted slices:
\[
\text{FAIL\_FRAC} = \sum_{s \in \text{slices}} \sum_{m \in \text{IQMs}} \mathbf{1}\{\text{IQM}_{m,s} \text{ is non-finite}\}
\]


% \section*{Meta Note}
% These metrics correspond to imaging physics QC literature and tools such as MRQy, MRIQC, ACR phantom procedures, and RadQy internal implementation.


%% REFERENCES %%%%%%%%%%%%%%%%%%%%%
\printbibliography[heading=bibintoc,title={References}]
\end{document}